\documentclass{article}
\usepackage[margin=2cm]{geometry}
\usepackage{amsmath}
\usepackage{graphicx}
\usepackage{booktabs}
\usepackage[utf8]{inputenc}
\usepackage{ragged2e}
\setlength{\emergencystretch}{3em}
\begin{document}
\sloppy

\noindent\hfill\includegraphics[height=1.8em]{assets/figure_header_logo.png}\par\noindent\hrule\vspace{4pt}

\subsection*{0.1Mass deformations}


In the following, we consider the coset with$n = 1$ i.e. a single vector multiplet. The coset
representative is expressed in terms of four scalars$\dot { \psi } ^ { i }, i = 0. 1, 2, 3$ via


\begin{equation}
L = \prod _ { i = 0 } ^ { 3 } e ^ { \phi ^ { i } K ^ { * } } \qquad ( 1 )
\end{equation}


where $K ^ { r }$ are the non compact generators of.${ \mathfrak { s O } } ( 4 ; 1 )$ see for details. Note that. $\phi ^ { 0 }$ is an
$S U ( \mathbf { 2 } ) . { R }$singlet, while the other three scalars form an$\phi ^ { r }$ $S U ( 2 ) . { R }$triplet. The scalar potential.
for this specific case can be obtained from and takes the following form


\begin{equation}
\begin{array} { r } { V ( \sigma, \phi ^ { i } ) = - g ^ { 2 } e ^ { 2 \sigma } + \frac { 1 } { 8 } m e ^ { - 6 \sigma } \bigg [ - 3 2 g e ^ { 4 \sigma } \cosh \phi ^ { 9 } \cosh \phi ^ { 1 } \cosh \phi ^ { 2 } \cosh \phi ^ { 3 } + 8 m \cosh ^ { 2 } \phi ^ { 0 } } \\ { \qquad }\end{array}
\end{equation}


The supersymmetric AdSvacuum is given by setting$\mathbf { 6 }$ $g = 3 m$ and setting all scalars to
vanish. The masses of the linearized scalar fluctuation around the AdS vacuum determine
the dimensions of the dual scalar operators in the SCFT via


\begin{equation}
m ^ { 2 } l ^ { 2 } = \Delta ( \Delta - 5 ) \qquad ( 3 )
\end{equation}


where is the curvature radius of the.$l$ $\mathbf { \hat { A } } \mathbf { \hat { C } } \mathbf { \hat { S } } . { i j }$ vacuum. For the scalars at hand, one finds


\begin{equation}
m _ { \sigma } ^ { 2 } l ^ { 2 } = - 6 ~ m _ { \phi } ^ { 2 } l ^ { 2 } = - 4 ~ m _ { \phi ^ { \prime } } ^ { 2 } l ^ { 2 } = - 6 ~, ~ ~ r = 1, 2, 3 ~
\end{equation}


Hence the dimensions of the dual operators are.


\begin{equation}
\Delta _ { \mathcal { O } _ { \varphi } } = 3, ~
\end{equation}


In these CFT operators were expressed in terms of free hypermultiplets (i.e. the singleton
sector). The case of${ \mathfrak { n } } = \mathbf { 1 }$ corresponds to having a single free hypermultiplet, consisting of
four real scalars and two symplectic Majorana spinors$q . { \cdot 1 } ^ { i }$ $\psi ^ { i }$ Here$i = 1, 2$ is the $S U ( 2 ) . { H }$
R-symmetry index and $A = 1, z$ is the $S U ( 2 )$ flavor symmetry index. The gauge invariant
operators appearing in are related to these fundamental fields as follows,.


\begin{equation}
\begin{array} { r l } { \mathcal { O } _ { \sigma } = ( q ^ { * } ) ^ { A } q _ { \mathrm { ~ A } } ^ { I }, \qquad \mathcal { O } _ { \phi ^ { 0 } } = \tilde { \psi } _ { f } \zeta ^ { I }, \qquad \mathcal { O } _ { \phi ^ { \prime } } = \left ( q ^ { * } \right ) ^ { A } ( \sigma ^ { \prime } ) ^ { B } q _ { \mathrm { ~ B } } ^ { I }, \quad r = 1, 2, 3 \qquad ( 6 ) } \\ { \mathrm { ~ o t h e ~ a d. ~ b e}}\end{array}
\end{equation}


Note that the first two operators correspond to mass terms for the scalars and fermions
respectively, in the hypermultiplet. The third operator is a triplet with respect to the.

\end{document}
